% Tiempos de cómputo de todos los análisis (tablas: reproducibility/figures_tables/make_timing_notes.py).

\snote{time}

This note reports the computing time of every analysis of the paper, problem by problem, with the purpose of showing what it costs to select and verify a subset with the method, how that cost compares with the standard identifiability analysis, profile likelihood of the full model, and where the computing time of the whole study was spent.

\subsection*{How the times were measured}
Each program records, for every problem, the elapsed time of its computation (field \texttt{segundos} of the result files of the reproducibility package). Several analyses use parallelism within a problem: the selection distributes the columns of $J$ and the stage-2 ordering over 16 to 64 processes; the nested validation distributes the finite-difference Jacobian of each fit of the full model; profile likelihood of the full model computes all points of a problem in parallel on up to 120 cores. Other analyses consist of many independent single-core tasks (the seeds of the synthetic experiment, the scenarios and seeds of the admissibility test of the other criteria, the points of the identifiability profiles); for them the tables give the sum over the tasks, which is processor time and is larger than the elapsed time. The problems with expensive simulations were computed on shared servers with 128 and 256 cores and the others on an 8-core workstation; the times on shared servers are upper bounds, because other computations ran at the same time. The complete method was timed in a run from scratch, without cached sensitivities (Supplementary Table~\ref{tab:tiempos} compares it with profile likelihood).

\subsection*{What the times show}
\emph{Selecting a subset is affordable.} The complete method took a median of 19~s per problem and less than ten minutes in 32 of the 34 problems; the slowest were Chen (33~min, 35~s per simulation), Smith (19~min) and Isensee (10~min). On one core the same computation took at most about three times longer (Supplementary Table~\ref{tab:timemethod}).

\emph{Verifying it is affordable too.} Determining the identifiability of the selected parameters (one fit of the reduced problem and the re-estimations at the profile points) took a median of 34~s per problem, and it was feasible in all 31 admissible problems, including those in which profile likelihood of the full model could not be computed (Supplementary Table~\ref{tab:timevalidation}). Profile likelihood of the full model, the standard identifiability analysis, took a median of 14~s in the 26 problems in which it was feasible (Supplementary Table~\ref{tab:timecomparison}), so on small problems it is as fast as our method; it could not be completed for the problems with more than 60 free parameters, in which the method and the identifiability check of the reduced problem still ran in minutes to hours.

\emph{Where the computing time went.} All analyses of the study together took about 168 hours of computing time (Supplementary Table~\ref{tab:timesummary}). Most of it was spent in analyses that re-estimate models many times and that are not part of the method: the variants of the method settings (50~h, because each variant reruns the complete method), the synthetic experiments (44~h over the three amplitudes), the identifiability profiles (26~h) and the tests of the subsets of the other criteria (25~h). Five problems with slow simulations or many calibration parameters, Isensee, Raimundez, Chen, Smith and Lang, account for 87\% of the total, whereas the median over the problems of the time of all analyses together is 21 minutes.

{\scriptsize\setlength{\tabcolsep}{3pt}
\begin{longtable}{@{}lrrlr@{}}
\caption{Summary of the computing time per analysis over the problems of Supplementary Tables~\ref{tab:timemethod}--\ref{tab:timecomparison}: number of problems, median and maximum time per problem and total time.}\label{tab:timesummary}\\
\toprule
Analysis & Problems & Median & Maximum (problem) & Total (h) \\
\midrule
\endfirsthead
\toprule
Analysis & Problems & Median & Maximum (problem) & Total (h) \\
\midrule
\endhead
Selection, parallel & 34 & 19 s & 32.6 min (Chen) & 1.3 \\
Selection, 1 core & 32 & 25 s & 28.3 min (Smith) & 1.3 \\
Settings (15 variants) & 28 & 2.6 min & 25.8 h (Raimundez) & 50.3 \\
Robustness, paired & 31 & 34 s & 22.1 min (Smith) & 1.4 \\
Robustness, independent & 31 & 27 s & 1.7 h (Smith) & 3.8 \\
Single parameter & 27 & 32 s & 1.6 h (Smith) & 3.2 \\
External validation & 29 & 7 s & 34.0 min (Smith) & 1.9 \\
Nested validation & 29 & $<$1 s & 11.0 min (Chen) & 0.4 \\
Synthetic, $\pm$20\%, 20 seeds & 29 & 2.8 min & 7.3 h (Chen) & 22.3 \\
Synthetic, $\pm$5\%, 10 seeds & 29 & 1.6 min & 2.8 h (Chen) & 8.6 \\
Synthetic, $\pm$50\%, 10 seeds & 29 & 2.6 min & 4.1 h (Chen) & 13.1 \\
True signal, 3 amplitudes & 29 & 31 s & 54.8 min (Chen) & 1.5 \\
Identifiability of $S$ & 31 & 34 s & 14.8 h (Chen) & 26.0 \\
Pivoted QR & 28 & $<$1 s & 1.9 min (Isensee) & 0.0 \\
Brun and Morris & 29 & 10 s & 29.6 min (Isensee) & 0.8 \\
Random subsets (20) & 29 & 4 s & 1.2 h (Isensee) & 1.6 \\
Calibration only & 28 & $<$1 s & 4.6 min (Isensee) & 0.1 \\
Criteria test, $\pm$5\% & 29 & 1.8 min & 13.8 h (Isensee) & 23.0 \\
Criteria test, $\pm$20\% & 29 & 25 s & 1.4 h (Smith) & 2.3 \\
Profile likelihood, full model & 26 & 14 s & 5.2 h (Isensee) & 6.2 \\
\midrule
All analyses (selection on one core excluded) & & & & 168 \\
\bottomrule
\end{longtable}

}

{\scriptsize\setlength{\tabcolsep}{3pt}
\begin{longtable}{@{}lrrrrrr@{}}
\caption{Computing time of the method and of the analyses of its behaviour, for every problem. Selection: complete method (sensitivities, stages 1 and 2, admissibility test and backward elimination) run from scratch, with the sensitivities in parallel or on one core; Settings: the 15 variants of Supplementary Note \emph{Sensitivity to the method settings} together; Robustness: seven perturbation sizes with paired or independent directions; Single parameter: Supplementary Note \emph{Single-parameter subsets}. Times recorded by each program for each problem: elapsed time of the computation of that problem with the parallelism described in the text; where an analysis consists of independent single-core tasks (seeds, scenarios, profile points) the time is their sum, i.e.\ processor time. --, not computed (problem not admissible, not applicable or not run). Large problems were computed on shared servers with 128 or 256 cores, the others on an 8-core workstation; times on shared servers may be inflated by resource contention.}\label{tab:timemethod}\\
\toprule
Problem & \rotatebox{70}{Selection, parallel} & \rotatebox{70}{Selection, 1 core} & \rotatebox{70}{Settings (15 variants)} & \rotatebox{70}{Robustness, paired} & \rotatebox{70}{Robustness, independent} & \rotatebox{70}{Single parameter} \\
\midrule
\endfirsthead
\toprule
Problem & \rotatebox{70}{Selection, parallel} & \rotatebox{70}{Selection, 1 core} & \rotatebox{70}{Settings (15 variants)} & \rotatebox{70}{Robustness, paired} & \rotatebox{70}{Robustness, independent} & \rotatebox{70}{Single parameter} \\
\midrule
\endhead
Alkan & 25 s & 46 s & -- & 47 s & 32 s & 1.0 h \\
Armistead & 7 s & 6 s & 56 s & 10 s & 6 s & 4 s \\
Bachmann & 10 s & 27 s & -- & 34 s & 16 s & -- \\
Beer & 3.3 min & 5.1 min & -- & -- & -- & -- \\
Bertozzi & 3 s & 4 s & 12 s & 4 s & 3 s & 2 s \\
Blasi & 4 s & 4 s & 1.3 min & 11 s & 9 s & 8 s \\
Boehm & 4 s & 3 s & 20 s & 11 s & 7 s & 6 s \\
Borghans & 18 s & 37 s & 2.6 min & 1.7 min & 1.1 min & 3.4 min \\
Brannmark & 27 s & 37 s & 3.5 min & 2.1 min & 1.6 min & 2.2 min \\
Bruno & 23 s & 34 s & 4.9 min & 22 s & 16 s & 8 s \\
Chen & 32.6 min & -- & -- & 7.0 min & 1.4 h & -- \\
Crauste & 36 s & 51 s & 3.3 min & -- & -- & -- \\
Elowitz & 12 s & 9 s & 2.2 min & 39 s & 28 s & 32 s \\
Fiedler & 17 s & 12 s & 1.9 min & 54 s & 38 s & 1.1 min \\
Fujita & 1.3 min & 1.3 min & 26.5 min & 9.6 min & 3.2 min & 3.0 min \\
Giordano & 12 s & 17 s & 1.0 min & 30 s & 20 s & 44 s \\
Isensee & 9.7 min & 23.6 min & -- & 13.5 min & 11.5 min & -- \\
Lang & 3.3 min & -- & 18.8 h & -- & -- & -- \\
Laske & 27 s & 28 s & 3.8 min & 4.6 min & 1.4 min & 30 s \\
Liu & 4 s & 6 s & 39 s & 46 s & 12 s & 12 s \\
Lucarelli & 35 s & 1.3 min & -- & 46 s & 11 s & 3.9 min \\
Okuonghae & 3 s & 6 s & 16 s & 23 s & 6 s & 15 s \\
Oliveira & 5 s & 11 s & 36 s & 49 s & 12 s & 18 s \\
Perelson & 2 s & 3 s & 2 s & 5 s & 1 s & $<$1 s \\
Rahman & 3 s & 4 s & 11 s & 14 s & 5 s & 5 s \\
Raia & 30 s & 49 s & 4.9 min & $<$1 s & 2.2 min & 2.3 min \\
Raimundez & 2.0 min & 5.1 min & 25.8 h & 2.6 min & 2.7 min & -- \\
Salazar-Cavazos & 1.1 min & 38 s & 5.0 min & 3 s & 1.5 min & 1.8 min \\
Schwen & 1.5 min & 2.5 min & 24.6 min & 10.7 min & 8.1 min & 14.3 min \\
Smith & 19.1 min & 28.3 min & 4.0 h & 22.1 min & 1.7 h & 1.6 h \\
Sneyd & 21 s & 18 s & 2.0 min & $<$1 s & 27 s & 36 s \\
Weber & 18 s & 23 s & 2.8 min & 32 s & 1.2 min & 2.8 min \\
Zhao & 8 s & 22 s & 2.5 min & 21 s & 17 s & 12 s \\
Zheng & 14 s & 17 s & 5.5 min & 22 s & 16 s & 21 s \\
\bottomrule
\end{longtable}

}

{\scriptsize\setlength{\tabcolsep}{3pt}
\begin{longtable}{@{}lrrrrrrr@{}}
\caption{Computing time of the validation and identifiability analyses, for every problem: external and nested validation on real data, synthetic experiment (20 seeds at $\pm$20\%, 10 seeds at $\pm$5\% and $\pm$50\%), true signal $y(\theta^\ast)-y(\theta_0)$ at the three amplitudes, and the identifiability of the selected parameters (fit of the reduced problem and all profile points). Times recorded by each program for each problem: elapsed time of the computation of that problem with the parallelism described in the text; where an analysis consists of independent single-core tasks (seeds, scenarios, profile points) the time is their sum, i.e.\ processor time. --, not computed (problem not admissible, not applicable or not run). Large problems were computed on shared servers with 128 or 256 cores, the others on an 8-core workstation; times on shared servers may be inflated by resource contention.}\label{tab:timevalidation}\\
\toprule
Problem & \rotatebox{70}{External validation} & \rotatebox{70}{Nested validation} & \rotatebox{70}{Synthetic, $\pm$20\%, 20 seeds} & \rotatebox{70}{Synthetic, $\pm$5\%, 10 seeds} & \rotatebox{70}{Synthetic, $\pm$50\%, 10 seeds} & \rotatebox{70}{True signal, 3 amplitudes} & \rotatebox{70}{Identifiability of $S$} \\
\midrule
\endfirsthead
\toprule
Problem & \rotatebox{70}{External validation} & \rotatebox{70}{Nested validation} & \rotatebox{70}{Synthetic, $\pm$20\%, 20 seeds} & \rotatebox{70}{Synthetic, $\pm$5\%, 10 seeds} & \rotatebox{70}{Synthetic, $\pm$50\%, 10 seeds} & \rotatebox{70}{True signal, 3 amplitudes} & \rotatebox{70}{Identifiability of $S$} \\
\midrule
\endhead
Alkan & 5.3 min & 6 s & 59.2 min & 23.6 min & 33.7 min & 1.5 min & 21.6 min \\
Armistead & $<$1 s & $<$1 s & 58 s & 10 s & 11 s & 26 s & $<$1 s \\
Bachmann & 4.8 min & 1.6 min & 1.6 h & 34.3 min & 53.5 min & 58 s & 35.2 min \\
Beer & -- & -- & -- & -- & -- & -- & -- \\
Bertozzi & -- & -- & -- & -- & -- & -- & 1 s \\
Blasi & -- & -- & -- & -- & -- & -- & 5 s \\
Boehm & $<$1 s & $<$1 s & 24 s & 12 s & 10 s & 27 s & 4 s \\
Borghans & 6 s & $<$1 s & 1.4 min & 1.1 min & 2.6 min & 26 s & 34 s \\
Brannmark & 8 s & $<$1 s & 2.8 min & 1.6 min & 1.6 min & 29 s & 58 s \\
Bruno & $<$1 s & $<$1 s & 18 s & 8 s & 9 s & 25 s & 19 s \\
Chen & 31.3 min & 11.0 min & 7.3 h & 2.8 h & 4.1 h & 54.8 min & 14.8 h \\
Crauste & -- & -- & -- & -- & -- & -- & -- \\
Elowitz & 6 s & $<$1 s & 1.4 min & 39 s & 42 s & 26 s & 22 s \\
Fiedler & 5 s & $<$1 s & 2.7 min & 1.1 min & 1.3 min & 43 s & 35 s \\
Fujita & 7 s & $<$1 s & 1.9 min & 1.0 min & 1.5 min & 30 s & 1.1 min \\
Giordano & 28 s & $<$1 s & 8.4 min & 5.0 min & 3.7 min & 43 s & 6 s \\
Isensee & 27.4 min & 10.6 min & 5.0 h & 1.7 h & 2.9 h & 13.2 min & 4.3 h \\
Lang & -- & -- & -- & -- & -- & -- & -- \\
Laske & 3 s & $<$1 s & 1.4 min & 42 s & 41 s & 43 s & 1.2 min \\
Liu & $<$1 s & $<$1 s & 27 s & 12 s & 10 s & 27 s & 7 s \\
Lucarelli & 3.2 min & 49 s & 51.7 min & 23.6 min & 29.2 min & 1.2 min & 54.9 min \\
Okuonghae & 7 s & $<$1 s & 1.8 min & 54 s & 49 s & 25 s & 4 s \\
Oliveira & 2 s & $<$1 s & 46 s & 20 s & 20 s & 28 s & 6 s \\
Perelson & $<$1 s & $<$1 s & 12 s & 6 s & 5 s & 22 s & $<$1 s \\
Rahman & $<$1 s & $<$1 s & 30 s & 10 s & 12 s & 25 s & 1 s \\
Raia & 18 s & $<$1 s & 3.2 min & 1.6 min & 2.6 min & 31 s & 1.4 min \\
Raimundez & 3 s & 1.2 min & 57.9 min & 37.0 min & 1.3 h & 2.7 min & 4.0 h \\
Salazar-Cavazos & 21 s & 2 s & 4.7 min & 2.1 min & 3.0 min & 2.5 min & 40 s \\
Schwen & 39 s & 2 s & 12.5 min & 7.4 min & 9.7 min & 34 s & 6.2 min \\
Smith & 34.0 min & 19 s & 4.0 h & 1.3 h & 1.7 h & 3.8 min & 48.1 min \\
Sneyd & 2 s & $<$1 s & 43 s & 19 s & 34 s & 29 s & 20 s \\
Weber & 1.1 min & $<$1 s & 13.4 min & 5.9 min & 9.2 min & 31 s & 37 s \\
Zhao & 22 s & $<$1 s & 5.7 min & 3.1 min & 4.0 min & 27 s & 9 s \\
Zheng & 1.3 min & $<$1 s & 27.4 min & 10.7 min & 21.3 min & 33 s & 9 s \\
\bottomrule
\end{longtable}

}

{\scriptsize\setlength{\tabcolsep}{3pt}
\begin{longtable}{@{}lrrrrrrr@{}}
\caption{Computing time of the comparison with other criteria, for every problem: selection and validation by pivoted QR, Brun and Morris, the random and calibration-only baselines, the admissibility and synthetic tests of the subsets of all criteria ($\pm$5\%, including the synthetic seeds, and $\pm$20\%), and profile likelihood of every biological parameter of the full model (all rounds). Times recorded by each program for each problem: elapsed time of the computation of that problem with the parallelism described in the text; where an analysis consists of independent single-core tasks (seeds, scenarios, profile points) the time is their sum, i.e.\ processor time. --, not computed (problem not admissible, not applicable or not run). Large problems were computed on shared servers with 128 or 256 cores, the others on an 8-core workstation; times on shared servers are upper bounds.}\label{tab:timecomparison}\\
\toprule
Problem & \rotatebox{70}{Pivoted QR} & \rotatebox{70}{Brun and Morris} & \rotatebox{70}{Random subsets (20)} & \rotatebox{70}{Calibration only} & \rotatebox{70}{Criteria test, $\pm$5\%} & \rotatebox{70}{Criteria test, $\pm$20\%} & \rotatebox{70}{Profile likelihood, full model} \\
\midrule
\endfirsthead
\toprule
Problem & \rotatebox{70}{Pivoted QR} & \rotatebox{70}{Brun and Morris} & \rotatebox{70}{Random subsets (20)} & \rotatebox{70}{Calibration only} & \rotatebox{70}{Criteria test, $\pm$5\%} & \rotatebox{70}{Criteria test, $\pm$20\%} & \rotatebox{70}{Profile likelihood, full model} \\
\midrule
\endhead
Alkan & 4 s & 1.9 min & 2.4 min & 5 s & 28.6 min & 50 s & 27.4 min \\
Armistead & $<$1 s & $<$1 s & $<$1 s & $<$1 s & 28 s & 14 s & 3 s \\
Bachmann & 1 s & 1.8 min & 10.1 min & 19 s & 1.5 h & 3.8 min & 1.2 h \\
Beer & -- & -- & -- & -- & -- & -- & -- \\
Bertozzi & $<$1 s & 2 s & $<$1 s & $<$1 s & 16 s & 6 s & $<$1 s \\
Blasi & -- & -- & -- & -- & -- & -- & -- \\
Boehm & $<$1 s & 3 s & $<$1 s & $<$1 s & 35 s & 10 s & $<$1 s \\
Borghans & $<$1 s & 6 s & 12 s & $<$1 s & 3.0 min & 1.3 min & 2.1 min \\
Brannmark & $<$1 s & 10 s & 5 s & $<$1 s & 4.3 min & 2.1 min & 13 s \\
Bruno & $<$1 s & 3 s & $<$1 s & $<$1 s & 39 s & 13 s & 2 s \\
Chen & -- & -- & -- & -- & -- & -- & -- \\
Crauste & -- & -- & -- & -- & -- & -- & -- \\
Elowitz & $<$1 s & 4 s & 2 s & $<$1 s & 1.1 min & 23 s & 10 s \\
Fiedler & $<$1 s & 6 s & 4 s & $<$1 s & 1.8 min & 34 s & $<$1 s \\
Fujita & $<$1 s & 9 s & 6 s & $<$1 s & 10.9 min & 2.9 min & 26 s \\
Giordano & $<$1 s & 12 s & $<$1 s & $<$1 s & 1.3 min & 23 s & 2.1 min \\
Isensee & 1.9 min & 29.6 min & 1.2 h & 4.6 min & 13.8 h & 25.5 min & 5.2 h \\
Lang & -- & -- & -- & -- & -- & -- & -- \\
Laske & $<$1 s & 11 s & 6 s & $<$1 s & 3.7 min & 1.2 min & 5 s \\
Liu & $<$1 s & 3 s & $<$1 s & $<$1 s & 35 s & 11 s & 1 s \\
Lucarelli & 3 s & 1.4 min & 52 s & $<$1 s & 11.1 min & 33 s & 10.9 min \\
Okuonghae & $<$1 s & 10 s & $<$1 s & $<$1 s & 34 s & 9 s & 23 s \\
Oliveira & $<$1 s & 2 s & 1 s & $<$1 s & 40 s & 12 s & 5 s \\
Perelson & $<$1 s & 2 s & $<$1 s & $<$1 s & 6 s & 1 s & $<$1 s \\
Rahman & $<$1 s & 3 s & $<$1 s & $<$1 s & 16 s & 4 s & 2 s \\
Raia & $<$1 s & 11 s & 11 s & $<$1 s & 4.0 min & 1.8 min & 3.5 min \\
Raimundez & 18 s & 7.5 min & 10 s & 2 s & 2.9 h & 5.4 min & -- \\
Salazar-Cavazos & 2 s & 12 s & 12 s & 2 s & 8.4 min & 1.6 min & 14 s \\
Schwen & $<$1 s & 28 s & 30 s & 2 s & 19.4 min & 6.5 min & 1.8 min \\
Smith & 5 s & 5.3 min & 8.0 min & 14 s & 3.1 h & 1.4 h & 14.2 min \\
Sneyd & $<$1 s & 3 s & 2 s & $<$1 s & 1.3 min & 25 s & 4 s \\
Weber & $<$1 s & 10 s & 4 s & $<$1 s & 2.4 min & 51 s & 2.0 min \\
Zhao & $<$1 s & 5 s & 2 s & $<$1 s & 44 s & 13 s & 20 s \\
Zheng & $<$1 s & 14 s & 2 s & $<$1 s & 1.0 min & 15 s & 4.8 min \\
\bottomrule
\end{longtable}

}
